\section{“Silicon Valley is LLM-pilled”：叙事、榜单与资源配置}

讲完 AMI 的组织选择之后，就能理解为什么访谈会把矛头指向“硅谷叙事”。这里的批评对象不是某个地理地点，而是一个会把问题、榜单和资源绑定在一起的产业系统。

访谈标题里的“硅谷被催眠了”，指的是整个 AI 行业被 LLM 叙事高度组织起来。谢赛宁并不是说 LLM 没价值，而是说一个叙事会定义 benchmark，benchmark 会定义 resource allocation，resource allocation 又会决定研究者能做什么。最后，很多有能力的人不是不想做世界模型、视频理解或物理智能，而是被组织目标分配到更贴近当前价值链的位置。

\subsection{价值链如何压缩研究空间}

访谈中的链条可以概括为：

\begin{center}
\begin{tabularx}{0.95\textwidth}{p{0.18\textwidth}p{0.28\textwidth}X}
\toprule
环节 & 表现 & 后果 \\
\midrule
叙事 & AGI、scaling law、LLM frontier & 定义什么问题“看起来重要”。 \\
榜单 & Chatbot Arena、数学、代码、通用能力排行 & 指挥团队优化可见指标。 \\
资源配置 & 算力、人力、发布节奏向榜单和产品集中 & 探索性研究缺氧。 \\
岗位分配 & 视频理解等问题被拆到 captioning 或产品支持环节 & 真正的 world-model-first 路线难以展开。 \\
\bottomrule
\end{tabularx}
\end{center}

\begin{warningbox}{被叙事催眠的风险}
当一个叙事过强时，研究者会把“对榜单有用”误认为“对智能本质有用”。这不是某个公司的问题，而是整个产业价值链对研究问题的筛选效应。
\end{warningbox}

\subsection{为什么逃出硅谷不是地理问题}

“逃出硅谷”不是简单的地理搬家。硅谷仍然有最密集的人才、资本和工程文化，AMI 未来也可能在硅谷设点。真正要逃出的，是单一 LLM 叙事、产品发布周期和榜单军备竞赛对问题定义的控制。

巴黎、纽约、蒙特利尔、新加坡的多地办公室在访谈里不只是行政安排，也对应“world model needs the world”的组织想象：不同地区、不同产业、不同学术网络、不同数据源共同参与。

\subsection{本章小结}

LLM-pilled 的核心问题不是 LLM 太强，而是叙事太单一。世界模型路线需要从现有价值链旁边开出空间，让视频、物理、机器人、工业过程和真实反馈重新成为 AI 研究的一等对象。

